\chapter{Budget}

The project budget is mainly a development and computation budget. No physical
prototype was manufactured. The main resources were student engineering time,
supervision time, local workstation use, and HPC GPU time.

\begin{table}[ht]
\centering
\caption{Estimated development budget. Monetary values are indicative and can
be adjusted to the official ETSETB costing criteria.}
\label{tab:budget}
\begin{tabularx}{\textwidth}{lrrX}
\toprule
Item & Quantity & Unit cost & Notes \\
\midrule
Student engineering work & 450 h & 0 EUR & Academic work; opportunity cost not billed. \\
Academic supervision & 25 h & 0 EUR & Included in degree supervision. \\
Cluster GPU compute, final Try 80 & 4 GPUs x 12 h & shared & RTX 2080 Ti-class run. \\
Cluster GPU compute, Tries 76--79 & 4 GPUs x 18 h & shared & Final-stage experiments before Try 80. \\
Exploratory compute & over 200 h & shared & Research iterations; not counted as deployable cost. \\
Storage and code hosting & existing & 0 EUR & Reused institutional/local infrastructure. \\
\bottomrule
\end{tabularx}
\end{table}

If electricity is costed separately, a representative final Try 80 run at
approximately \SI{1.5}{kW} for \SI{12}{h} consumes about \SI{18}{kWh}. At
\SI{0.20}{EUR/kWh}, this corresponds to approximately 3.6 EUR of electricity
for the final run. This value is small compared with personnel time and also
small compared with repeatedly running full ray-tracing simulations for new
deployments.
